Three findings stand out.
\textbf{(i) With exact anchors, the bias disappears and the MSE drops by orders of magnitude at high noise.}
With TSI at $\sigma_t=70$, the relative MSE falls from $2.4\cdot10^6$ (iDEM, $K=10$) to $5.0\cdot10^3$ (AEM, $K=10$), a $475\times$ reduction. The gain is still $29\times$ at $K=500$.
Even AEM with $K=0$ (pure target score matching on the anchor) beats iDEM with $K=500$ at high noise. Conversely, at low noise ($\sigma_t\le0.25$) the two coincide, as expected: there the proposals already sample the posterior well.
With DSI, iDEM's relative squared bias stays around $7$--$8\cdot10^{-2}$ at $K=10$ for all $\sigma_t\ge2.8$. AEM's is indistinguishable from zero, as \cref{thm:unbiased} predicts.
\textbf{(ii) TSI is the wrong test function at high noise for both methods.} Its relative variance grows like $\sigma_t^2$ because the signal $\nabla\log p_t=O(\sigma_t^{-1})$ shrinks while $\mathrm{Var}[\nabla\mathcal{E}]$ does not. DSI is accurate there but noisy at low noise. This motivates the precision-weighted combination (\cref{sec:method}), which all methods use in \cref{sec:exp-training}. For a unit Gaussian component the combination has zero variance.
\textbf{(iii) Approximate anchors behave as \cref{thm:contraction} predicts.} With tempered anchors (component std $2$ instead of $1$), the bias of AEM at $\sigma_t=2.8$ shrinks geometrically in $K$ ($8.9\cdot10^{-2}$, $2.1\cdot10^{-2}$, $2.5\cdot10^{-3}$, $5\cdot10^{-4}$ for $K=0,10,100,500$). It is already below iDEM's bias at $K\ge10$.
With anchors drawn from a target that lacks half of the modes, the proposals at intermediate noise ($\sigma_t\approx14$) rarely reach the missing modes, so $\epsilon_K$ is small. AEM's bias then \emph{exceeds} iDEM's ($0.16$ vs.\ $0.05$ at $K=10$). This is the buffer-reliance regime we return to in \cref{sec:exp-ablation}.
